\begin{figure}[htbp]
\centering
\begin{tikzpicture}
\begin{axis}[
  width=0.82\textwidth, height=6.0cm,
  xmin=0, xmax=250, ymin=0, ymax=22,
  xlabel={Cost of one die-to-die hop in the model (ns)}, ylabel={Round trip, other CCD $\div$ same CCD},
  legend columns=2, legend style={font=\scriptsize,at={(0.5,-0.34)},anchor=north}, legend cell align=left,
  grid=major, grid style={draftgray!20}]
\addplot[draftblue,very thick,mark=*,mark size=1.3pt] coordinates {(0,1.48) (9,2.214) (18,2.949) (20,3.112) (29,3.847) (37,4.5) (40,4.745) (60,6.378) (100,9.643) (150,13.72) (241,21.15)};
\addlegendentry{Model (ping-pong on one cache line)}
\addplot[mmred,thick,dashed] coordinates {(0,4.52) (250,4.52)};
\addlegendentry{Measured on Milan-X: 265 $\div$ 58.6\,ns (Reguly)}
\addplot[draftgray,thick,dotted] coordinates {(9,0) (9,11)};
\addlegendentry{AMD's CCD--IOD round trip, 9\,ns}
\addplot[draftgray,thick,dashdotted] coordinates {(18,0) (18,11)};
\addlegendentry{Two such round trips, 18\,ns}
\end{axis}
\end{tikzpicture}
\caption[S4: cross-CCD to same-CCD ping-pong ratio against hop cost]{S4: the model's ratio of cross-CCD to same-CCD ping-pong time as the cost of a hop grows (planar Zen~3 CCDs at 4\,GHz). Only the ratio is fitted to Reguly's measurement \cite{regulyComparativeEvaluationBandwidthbound2023}, because Reguly's test times one message while the model's round trip contains two messages and two write upgrades. One transfer between two different CCDs, a request and its reply, crosses the CCD--IOD link four times, two of AMD's 9\,ns round trips. With a 9\,ns hop the model gives 2.2$\times$ and with 18\,ns 2.9$\times$; matching the measured 4.5$\times$ needs a hop of about 37\,ns. That value is fitted to the ratio only: the model's exchanges are about 2.4 times faster than the measured ones, so 37\,ns is not the real cost of a transfer, which \cref{sec:cross-ccd} puts at about 206\,ns on one cloud instance.}
\label{fig:sim-pingpong}
\end{figure}
